\begin{figure*}[htb]
\centering

\tikzset{
    >=Stealth,
    font=\small,
    panelbox/.style={
        rounded corners,
        draw=black!60,
        very thick,
        inner sep=4mm,
        fill=white,
    },
    featurerow/.style={
        draw=black!60,
        rounded corners,
        very thin,
        fill=black!1,
        minimum height=4mm,
        anchor=west
    },
    heatrow/.style={
        draw=black!60,
        rounded corners,
        very thin,
        minimum height=4mm,
        anchor=west
    },
    mininode/.style={
        draw=black,
        thick,
        rounded corners,
        fill=black!2,
        minimum height=0.6cm,
        inner sep=2mm,
        align=center,
        font=\scriptsize
    },
    stage/.style={
        draw=black,
        thick,
        rounded corners,
        minimum width=3cm,
        minimum height=0.9cm,
        align=center,
        fill=black!3
    },
    miniStage/.style={
        draw=black,
        thick,
        rounded corners,
        fill=black!2,
        minimum width=1.0cm,
        minimum height=0.7cm,
        align=center,
        font=\scriptsize
    },
}

\begin{tikzpicture}[
    node distance=10mm and 20mm
]

% =====================================================
% LEFT PANEL: Feature-space explanation
% =====================================================

% Block containing features + attribution bars
\node[draw=black, thick, rounded corners,
      fill=white,
      minimum width=5.2cm,
      minimum height=2.8cm,
      inner sep=2mm,
      align=left] (fs_block) {};

% Populate rows inside fs_block using local coordinates
\begin{scope}[shift={(fs_block.north west)}]
    % column headers
    \node[anchor=south west,font=\scriptsize]
        at ([xshift=2mm,yshift=0mm]0,0) {feature $x_i$};
    \node[anchor=south west,font=\scriptsize]
        at ([xshift=3.4cm,yshift=0mm]0,0) {saliency $a_i$};

    % row 1
    \node[featurerow, minimum width=3.0cm, fill=black!1]
        (fx1) at ([xshift=2mm,yshift=-4mm]0,0)
        {\scriptsize \texttt{SectionEntropy}};
    \node[heatrow, minimum width=1.7cm, fill=black!70]
        (fa1) at ([xshift=3.4cm,yshift=-4mm]0,0) {};

    % row 2
    \node[featurerow, minimum width=3.0cm, fill=black!1]
        (fx2) at ([xshift=2mm,yshift=-10mm]0,0)
        {\scriptsize \texttt{NumberOfImports}};
    \node[heatrow, minimum width=1.2cm, fill=black!45]
        (fa2) at ([xshift=3.4cm,yshift=-10mm]0,0) {};

    % row 3
    \node[featurerow, minimum width=3.0cm, fill=black!1]
        (fx3) at ([xshift=2mm,yshift=-16mm]0,0)
        {\scriptsize \texttt{MedianStringLength}};
    \node[heatrow, minimum width=0.5cm, fill=black!20]
        (fa3) at ([xshift=3.4cm,yshift=-16mm]0,0) {};

    % row 4
    \node[featurerow, minimum width=3.0cm, fill=black!1]
        (fx4) at ([xshift=2mm,yshift=-22mm]0,0)
        {\scriptsize \texttt{HeaderPermissions}};
    \node[heatrow, minimum width=1.4cm, fill=black!55]
        (fa4) at ([xshift=3.4cm,yshift=-22mm]0,0) {};
\end{scope}

% Create some top padding so the panel label can sit above fs_block
\coordinate (fs_pad) at ($(fs_block.north)+(0,12mm)$);

% Put fs_block in a rounded "panelbox"
\begin{scope}[on background layer]
\node[panelbox,
      yshift=1.25cm,
      minimum width=8cm,
      minimum height=8cm,
      fit=(fs_pad) (fs_block),
      inner sep=6mm] (fs_panel) {};
\end{scope}

% Now label that panel the same way you did in your other figures
\node[font=\footnotesize,
      anchor=north,
      align=center]
    at ($(fs_panel.north)+(0,-2mm)$)
    {feature-space explanation\\
     $a \in \mathcal{A} \subseteq \mathbb{R}^{|x|}$, e.g., $a \in \mathbb{R}^{2^{10}}$};


\begin{scope}[shift={($(fs_panel.north west)+(12mm,-22mm)$)}]

    % left stack of inputs
    \node[miniStage] (xnode) {$x$};
    \node[miniStage, below=1mm of xnode] (fs_phinode) {$\phi$};

    % helper coordinate at vertical midpoint between x and phi
    \path (xnode) -- (fs_phinode) node[midway] (xphi_mid) {};

    % explanation algorithm box, centered vertically on that midpoint
    \node[stage,
          minimum width=2.0cm,
          align=center,
          right=8mm of xphi_mid] (fs_epsnode)
    {explanation alg.\\$\epsilon_{\phi}: \mathcal{X}\mapsto\mathcal{A}$};

    % final feature-space explanation box
    \node[stage,
          minimum width=2.0cm,
          align=center,
          right=3mm of fs_epsnode] (anode)
    {explanation\\$a \in \mathcal{A} \subseteq \mathbb{R}^{|x|}$};

    % arrows
    \draw[->,thick] (xnode.east) -- (fs_epsnode.west);
    \draw[->,thick] (fs_phinode.east) -- (fs_epsnode.west);
    \draw[->,thick] (fs_epsnode.east) -- (anode.west);

\end{scope}

% =====================================================
% RIGHT PANEL: Problem-space explanation
% =====================================================

% Block containing bytes / spans + attribution bars
\node[draw=black, thick, rounded corners,
      fill=white,
      minimum width=5.6cm,
      minimum height=2.8cm,
      inner sep=2mm,
      align=left,
      right=30mm of fs_block] (ps_block) {};

% Populate rows inside ps_block using local coordinates
\begin{scope}[shift={(ps_block.north west)}]
    % column headers
    \node[anchor=south west,font=\scriptsize]
        at ([xshift=2mm,yshift=0mm]0,0) {element(s) $z_i$};
    \node[anchor=south west,font=\scriptsize]
        at ([xshift=3.9cm,yshift=0mm]0,0) {saliency $b_i$};

    % row 1
    \node[featurerow, minimum width=3.5cm, fill=black!1]
        (pz1) at ([xshift=2mm,yshift=-4mm]0,0)
        {\scriptsize \texttt{4D 5A}};
    \node[heatrow, minimum width=0.2cm, fill=black!10]
        (pb1) at ([xshift=3.9cm,yshift=-4mm]0,0) {};

    % row 2
    \node[featurerow, minimum width=3.5cm, fill=black!1]
        (pz2) at ([xshift=2mm,yshift=-10mm]0,0)
        {\scriptsize \texttt{.text[0x401000:40]}};
    \node[heatrow, minimum width=1.6cm, fill=black!70]
        (pb2) at ([xshift=3.9cm,yshift=-10mm]0,0) {};

    % row 3
    \node[featurerow, minimum width=3.5cm, fill=black!1]
        (pz3) at ([xshift=2mm,yshift=-16mm]0,0)
        {\scriptsize \texttt{.data[0x402000:20]}};
    \node[heatrow, minimum width=0.7cm, fill=black!30]
        (pb3) at ([xshift=3.9cm,yshift=-16mm]0,0) {};

    % row 4
    \node[featurerow, minimum width=3.5cm, fill=black!1]
        (pz4) at ([xshift=2mm,yshift=-22mm]0,0)
        {\scriptsize \texttt{overlay}};
    \node[heatrow, minimum width=1.4cm, fill=black!55]
        (pb4) at ([xshift=3.9cm,yshift=-22mm]0,0) {};
\end{scope}

% Create top padding for the right panel
\coordinate (ps_pad) at ($(ps_block.north)+(0,12mm)$);

% Put ps_block in its rounded "panelbox"
\begin{scope}[on background layer]
\node[panelbox,
      yshift=1.25cm,
      fit=(ps_pad) (ps_block),
      minimum width=8cm,
      minimum height=8cm,
      inner sep=6mm] (ps_panel) {};
\end{scope}

% Label this panel at its north edge
\node[font=\footnotesize,
      anchor=north,
      align=center]
    at ($(ps_panel.north)+(0,-2mm)$)
    {problem-space explanation
     \\$b \in \mathcal{B} \subseteq \mathbb{R}^{|z|}$, e.g., $b \in \mathbb{R}^{2^{30}}$
     };

% --- flow diagram (inputs -> epsilon -> b), positioned inside ps_panel ---
\begin{scope}[shift={($(ps_panel.north west)+(12mm,-14mm)$)}]

    % left stack of inputs
    \node[miniStage] (znode) {$z$};
    \node[miniStage, below=1mm of znode] (sigmanode) {$\sigma$};
    \node[miniStage, below=1mm of sigmanode] (phinode) {$\phi$};

    % explanation algorithm box, constrained width so it wraps
    \node[stage,
          % text width=2.0cm,
          minimum width=2.0cm,
          align=center,
          right=4mm of sigmanode] (epsnode)
    {explanation alg.\\$\epsilon_{\sigma, \phi}: \mathcal{Z}\mapsto\mathcal{B}$};

    % final problem-space explanation box
    \node[stage,
          % text width=2.0cm,
          minimum width=2.0cm,
          align=center,
          right=3mm of epsnode] (bnode)
    {explanation\\$b \in \mathcal{B} \subseteq \mathbb{R}^{|z|}$};

    % arrows
    \draw[->,thick] (znode.east) -- (epsnode.west);
    \draw[->,thick] (sigmanode.east) -- (epsnode.west);
    \draw[->,thick] (phinode.east) -- (epsnode.west);
    \draw[->,thick] (epsnode.east) -- (bnode.west);

\end{scope}

\end{tikzpicture}

\caption{Explanatory Targets. Left: the feature-space explanation process, where the explanation algorithm $\epsilon_{\phi}$ provides saliency $a$ over the feature-space; and right: the problem-space explanation process, where the explanation algorithm $\epsilon_{\sigma,\phi}$ yield saliency $b$ over the problem-space.}
\label{fig:feature-vs-problem-explanations}
\end{figure*}
